\ifdefined\gkver\else\def\gkver{student}\fi
\documentclass[\gkver]{gaokaozhenti}
\begin{document}

\chapter{1986年上海}

\begin{enumerate}
\item
%% number: 1
%% paperName: 1986年上海高考第 1 题 2 分
%% typeId: 3
%% type: 填空题
%% chapter: 光学
%% point: 光的折射
%% method: 无
%% score: 2
%% degree: 850
%% duplicateId: 0
%% body:
白光从空气中进入玻璃三棱镜时会产生色散现象。\tkanswer{紫}色光向棱镜底边偏折最大；\tkanswer{红}色光在三棱镜中传播速率最大。

%% answer: 紫；红
\jdanswer{
紫；红
}

%% memo:
\memoanswer{本题原卷及材料中未提供详解，待补充。}

\item
%% number: 2
%% paperName: 1986年上海高考第 2 题 2 分
%% typeId: 3
%% type: 填空题
%% chapter: 交流电
%% point: LC振荡回路
%% method: 无
%% score: 2
%% degree: 850
%% duplicateId: 0
%% body:
LC 振荡电路的振荡频率 $f=$\tkanswer{$\frac{1}{2\pi\sqrt{LC}}$}。

%% answer: \frac{1}{2\pi\sqrt{LC}}
\jdanswer{
$\frac{1}{2\pi\sqrt{LC}}$
}

%% memo:
\memoanswer{本题原卷及材料中未提供详解，待补充。}

\item
%% number: 3
%% paperName: 1986年上海高考第 3 题 2 分
%% typeId: 3
%% type: 填空题
%% chapter: 交流电
%% point: 交流电
%% method: 无
%% score: 2
%% degree: 850
%% duplicateId: 0
%% body:
瞬时值为 $u=50\sqrt{2}\sin(100\pi t)$ 的交流电压，加在阻值 $R=4\UO$ 的电阻上，则流过电阻的电流有效值是\tkanswer{$12.5\UA$}；频率是\tkanswer{$50\UHz$}。

%% answer: 12.5 A；50 Hz
\jdanswer{
$12.5\UA$；$50\UHz$
}

%% memo:
\memoanswer{本题原卷及材料中未提供详解，待补充。}

\item
%% number: 4
%% paperName: 1986年上海高考第 4 题 2 分
%% typeId: 3
%% type: 填空题
%% chapter: 磁场
%% point: 洛伦兹力
%% method: 无
%% score: 2
%% degree: 850
%% duplicateId: 0
%% body:
一个电子在匀强磁场中运动而不受到磁场力的作用，则电子运动的方向是\tkanswer{平行于磁场方向}。

%% answer: 平行于磁场方向
\jdanswer{
平行于磁场方向
}

%% memo:
\memoanswer{本题原卷及材料中未提供详解，待补充。}

\item
%% number: 5
%% paperName: 1986年上海高考第 5 题 2 分
%% typeId: 3
%% type: 填空题
%% chapter: 机械振动
%% point: 单摆
%% method: 无
%% score: 2
%% degree: 650
%% duplicateId: 0
%% body:
单摆作简谐振动，频率为 $0.5\UHz$，振幅为 $4\Ucm$。当摆球具有最大正向速度的瞬时开始计时，在右图中画出振动图象。
\onepicture[align=right, w=6cm]{\includesvg{figs/05}}

%% answer: 图象为周期 2 s、振幅 4 cm、初相 0 的正弦曲线
\jdanswer{
图象应是周期是 $2\Us$、振幅是 $4\Ucm$、初位相是 $0$ 的正弦曲线。

\onepicture[w=6cm]{\includesvg{figs/05_an}}
}

%% memo:
\memoanswer{本题原卷及材料中未提供详解，待补充。}

\item
%% number: 6
%% paperName: 1986年上海高考第 6 题 2 分
%% typeId: 3
%% type: 填空题
%% chapter: 原子和原子核
%% point: 玻尔模型
%% method: 无
%% score: 2
%% degree: 850
%% duplicateId: 0
%% body:
在氢光谱巴耳末公式 $\frac{1}{\lambda}=R\left(\frac{1}{2^{2}}-\frac{1}{n^{2}}\right)$ 中，$R$ 为里德伯恒量。式中 $n$ 必须满足的条件是\tkanswer{大于 2 的正整数}。

%% answer: 大于 2 的正整数
\jdanswer{
大于 2 的正整数
}

%% memo:
\memoanswer{本题原卷及材料中未提供详解，待补充。}

\item
%% number: 7
%% paperName: 1986年上海高考第 7 题 2 分
%% typeId: 3
%% type: 填空题
%% chapter: 原子和原子核
%% point: 天然放射性
%% method: 无
%% score: 2
%% degree: 650
%% duplicateId: 0
%% body:
一个原子核经一次 $\alpha$ 衰变和一次 $\beta$ 衰变后，成为一新原子核。则新核与原来的核相比，质子少了\tkanswer{1}个；中子少了\tkanswer{3}个。

%% answer: 1；3
\jdanswer{
1；3
}

%% memo:
\memoanswer{本题原卷及材料中未提供详解，待补充。}

\item
%% number: 8
%% paperName: 1986年上海高考第 8 题 2 分
%% typeId: 3
%% type: 填空题
%% chapter: 运动和力
%% point: 动量守恒定律
%% method: 无
%% score: 2
%% degree: 650
%% duplicateId: 0
%% body:
甲、乙两溜冰者，质量分别为 $50\Ukg$ 和 $52\Ukg$，甲手里拿着一质量为 $2\Ukg$ 的球，两人均以 $2\Ums$ 的速度在冰面上相向滑行。甲将球抛给乙，乙再将球抛给甲，这样抛接若干次后，乙的速度变为零，则甲的速度是\tkanswer{$0$}。

%% answer: 0
\jdanswer{
0
}

%% memo:
\memoanswer{本题原卷及材料中未提供详解，待补充。}

\item
%% number: 9
%% paperName: 1986年上海高考第 9 题 3 分
%% typeId: 1
%% type: 单选题
%% chapter: 直线运动
%% point: 加速度
%% method: 无
%% score: 3
%% degree: 850
%% duplicateId: 0
%% body:
关于速度和加速度的关系，下列说法中正确的是\xzanswer{B}
\fourchoices[answer=B]
{速度变化得越多，加速度就越大}
{速度变化得越快，加速度就越大}
{加速度方向保持不变，速度方向也保持不变}
{加速度大小不断变小，速度大小也不断变小}

%% answer: B
%% memo:
\memoanswer{
“速度变化得越多”是指 $\Delta v$ 越大，若所用时间 $t$ 也很大，则 $\frac{\Delta v}{\Delta t}$ 就不一定大，故 A 错。“速度变化得越快”是指速度的变化率 $\frac{\Delta v}{\Delta t}$ 越大，即加速度 $a$ 越大，B 正确。加速度方向保持不变，速度方向可能变，也可能不变，当物体做减速直线运动时，$v=0$ 以后就反向运动，如竖直向上抛出的物体，故 C 错。物体在运动过程中，若加速度方向与速度方向相同，尽管加速度在变小，但物体仍在加速，直到加速度 $a=0$ 时，速度就达到最大了，故 D 错。
}

\item
%% number: 10
%% paperName: 1986年上海高考第 10 题 3 分
%% typeId: 1
%% type: 单选题
%% chapter: 运动和力
%% point: 超重失重
%% method: 无
%% score: 3
%% degree: 850
%% duplicateId: 0
%% body:
一质量为 $m$ 千克的物体挂在弹簧秤下，手持弹簧秤的上端加速上提，弹簧秤的读数为 $p$ 牛顿，则上提的加速度是\xzanswer{C}
\fourchoices[answer=C]
{$\frac{p}{m}$}
{$g$}
{$\frac{p}{m}-g$}
{$\frac{p}{m}+g$}

%% answer: C
%% memo:
\memoanswer{本题原卷及材料中未提供详解，待补充。}

\item
%% number: 11
%% paperName: 1986年上海高考第 11 题 3 分
%% typeId: 1
%% type: 单选题
%% chapter: 气体性质
%% point: 物体的内能
%% method: 无
%% score: 3
%% degree: 650
%% duplicateId: 0
%% body:
关于物体的内能，以下说法中正确的是\xzanswer{D}
\fourchoices[answer=D]
{不同的物体，若温度相等，则内能也相等}
{物体速度增大，则分子动能增大，内能也增大}
{晶体熔解时，温度不变，则内能也不变}
{对物体做功，或向物体传热，都可能改变物体的内能}

%% answer: D
%% memo:
\memoanswer{本题原卷及材料中未提供详解，待补充。}

\item
%% number: 12
%% paperName: 1986年上海高考第 12 题 3 分
%% typeId: 1
%% type: 单选题
%% chapter: 电磁感应
%% point: 楞次定律推广
%% method: 无
%% score: 3
%% degree: 650
%% duplicateId: 0
%% body:
一无限长直导线通有电流 $I$，有一矩形线圈与其共面，如图所示。当电流 $I$ 减小时，矩形线圈将\xzanswer{A}
\onepicture[w=2.2cm]{\includesvg{figs/12}}
\fourchoices[answer=A]
{向左平动}
{向右平动}
{静止不动}
{发生转动}

%% answer: A
%% memo:
\memoanswer{本题原卷及材料中未提供详解，待补充。}

\item
%% number: 13
%% paperName: 1986年上海高考第 13 题 3 分
%% typeId: 1
%% type: 单选题
%% chapter: 光学
%% point: 光的干涉
%% method: 无
%% score: 3
%% degree: 650
%% duplicateId: 0
%% body:
能产生干涉现象的两束光是\xzanswer{B}
\fourchoices[answer=B]
{频率相同、振幅相同的两束光}
{频率相同、位相差恒定的两束光}
{两只完全相同的灯泡发出的光}
{同一光源的两个发光部分发出的光}

%% answer: B
%% memo:
\memoanswer{本题原卷及材料中未提供详解，待补充。}

\item
%% number: 14
%% paperName: 1986年上海高考第 14 题 3 分
%% typeId: 1
%% type: 单选题
%% chapter: 光学
%% point: 光电效应
%% method: 无
%% score: 3
%% degree: 650
%% duplicateId: 0
%% body:
频率为 $\nu$ 的光照射某种金属材料，产生光电子的最大初动能为 $E_{k}$。若以频率为 $2\nu$ 的光照射同一金属材料，则光电子的最大初动能是\xzanswer{B}
\fourchoices[answer=B]
{$2E_{k}$}
{$E_{k}+h\nu$}
{$E_{k}-h\nu$}
{$E_{k}+2h\nu$}

%% answer: B
%% memo:
\memoanswer{本题原卷及材料中未提供详解，待补充。}

\item
%% number: 15
%% paperName: 1986年上海高考第 15 题 3 分
%% typeId: 1
%% type: 单选题
%% chapter: 原子和原子核
%% point: 玻尔模型
%% method: 无
%% score: 3
%% degree: 650
%% duplicateId: 0
%% body:
光子能量为 $E$ 的一束光，照射容器中的氢气。氢原子吸收光子后，能发出频率分别为 $\nu_{1}$、$\nu_{2}$、$\nu_{3}$ 的三种光，且 $\nu_{1}<\nu_{2}<\nu_{3}$。则入射光束的光子能量是\xzanswer{C}
\fourchoices[answer=C]
{$h\nu_{1}$}
{$h\nu_{2}$}
{$h\nu_{3}$}
{$h(\nu_{1}+\nu_{2}+\nu_{3})$}

%% answer: C
%% memo:
\memoanswer{本题原卷及材料中未提供详解，待补充。}

\item
%% number: 16
%% paperName: 1986年上海高考第 16 题 3 分
%% typeId: 1
%% type: 单选题
%% chapter: 机械波
%% point: 波长频率波速
%% method: 无
%% score: 3
%% degree: 650
%% duplicateId: 0
%% body:
有一频率为 $0.5\UHz$ 的简谐横波，沿 $x$ 轴正方向传播。在传播方向上有相距为 $2\Um$ 的两个质点，它们相继达到正向位移最大值的时间差为 $0.5\Us$，则波长是\xzanswer{D}
\fourchoices[answer=D]
{$1\Um$}
{$2\Um$}
{$4\Um$}
{$8\Um$}

%% answer: D
%% memo:
\memoanswer{本题原卷及材料中未提供详解，待补充。}

\item
%% number: 17
%% paperName: 1986年上海高考第 17 题 3 分
%% typeId: 2
%% type: 多选题
%% chapter: 其他
%% point: 物理学史
%% method: 无
%% score: 3
%% degree: 850
%% duplicateId: 0
%% body:
下列叙述中，符合物理学史事实的有\xzanswer{ABD}
\fourchoices[answer=ABD]
{电子是由汤姆生发现的}
{“相对论”是由爱因斯坦创立的}
{原子的核式结构学说是由玻尔提出的}
{理论上预言了电磁波存在的是麦克斯韦}

%% answer: ABD
%% memo:
\memoanswer{本题原卷及材料中未提供详解，待补充。}

\item
%% number: 18
%% paperName: 1986年上海高考第 18 题 3 分
%% typeId: 2
%% type: 多选题
%% chapter: 气体性质
%% point: 理想气体状态方程
%% method: 无
%% score: 3
%% degree: 650
%% duplicateId: 0
%% body:
对于理想气体，下列说法中正确的是\xzanswer{AD}
\fourchoices[answer=AD]
{在质量和温度不变的情况下，压强与体积成反比}
{在质量和压强不变的情况下，每升高 $1\UdC$，体积就增大 $\frac{1}{273}$}
{只有在温度不太低、压强不太大的情况下，才能服从理想气体状态方程}
{分子间无相互作用的势能，分子无规则运动的动能的总和就是内能}

%% answer: AD
%% memo:
\memoanswer{本题原卷及材料中未提供详解，待补充。}

\item
%% number: 19
%% paperName: 1986年上海高考第 19 题 3 分
%% typeId: 2
%% type: 多选题
%% chapter: 电场
%% point: 电势能电势差电场力做功
%% method: 无
%% score: 3
%% degree: 650
%% duplicateId: 0
%% body:
一正电荷在电场中沿某一电场线从 A 点移到 B 点，在此过程中有可能的是\xzanswer{ABD}
\onepicture[w=5cm]{\includesvg{figs/19}}
\fourchoices[answer=ABD]
{电场力的大小不断变大}
{电场力的大小保持不变}
{电荷克服电场力做功}
{电荷的电势能不断变小}

%% answer: ABD
%% memo:
\memoanswer{本题原卷及材料中未提供详解，待补充。}

\item
%% number: 20
%% paperName: 1986年上海高考第 20 题 3 分
%% typeId: 2
%% type: 多选题
%% chapter: 运动和力
%% point: 动量和冲量
%% method: 无
%% score: 3
%% degree: 650
%% duplicateId: 0
%% body:
一物体沿光滑斜面下滑，在此过程中\xzanswer{AC}
\fourchoices[answer=AC]
{斜面对物体的弹力做功为零}
{斜面对物体的弹力的冲量为零}
{物体动能的增量等于重力所做的功}
{物体动量的增量等于重力的冲量}

%% answer: AC
%% memo:
\memoanswer{本题原卷及材料中未提供详解，待补充。}

\item
%% number: 21
%% paperName: 1986年上海高考第 21 题 5 分
%% typeId: 4
%% type: 实验题
%% chapter: 力物体的平衡
%% point: 力矩盘
%% method: 无
%% score: 5
%% degree: 650
%% duplicateId: 0
%% body:
某学生在做“有固定转动轴物体的平衡条件”实验时，他把力矩盘调节到平衡，如图所示。盘上各圆的半径分别是 $0.1\Um$，$0.2\Um$，$0.3\Um$，$0.4\Um$，$0.5\Um$。每个钩码的质量均为 $0.1\Ukg$。若规定逆时针力矩为正，顺时针力矩为负。则 $F_{1}$ 的力矩是\tkanswer{$-0.12\Ukgf\cdot\Um$}；$F_{2}$ 的力矩是\tkanswer{$0.06\Ukgf\cdot\Um$}。根据平衡条件，测力计与圆盘连线上的拉力 $T$ 应该是\tkanswer{$0.3\Ukgf$}，但该学生发现测力计的读数与该值有偏差，除摩擦等原因外，从所示的两图中可看出引起误差的原因是：\tkanswer[4cm]{力矩盘盘面不在竖直平面里，测力计和拉线不在同一直线上}。
\twopicture[align=right, showlabel=false, h=5cm]{figs/21a.png}{figs/21b.png}

%% answer: -0.12 kgf·m，0.06 kgf·m，0.3 kgf，力矩盘盘面不在竖直平面里，测力计和拉线不在同一直线上
\jdanswer{
$-0.12\Ukgf\cdot\Um$，$0.06\Ukgf\cdot\Um$，$0.3\Ukgf$，力矩盘盘面不在竖直平面里，测力计和拉线不在同一直线上
}

%% memo:
\memoanswer{本题原卷及材料中未提供详解，待补充。}

\item
%% number: 22
%% paperName: 1986年上海高考第 22 题 6 分
%% typeId: 4
%% type: 实验题
%% chapter: 电路
%% point: 测电动势与内阻
%% method: 无
%% score: 6
%% degree: 650
%% duplicateId: 0
%% body:
给你一个伏特表，一个电阻箱，电键及导线等，如何根据全电路欧姆定律测出一节旧干电池的电动势和内电阻。
\begin{enumerate}
	\item
	在右方空白处画出实验电路图。

	\item
	实验过程中，将电阻箱拨到 $45\UO$ 时，伏特表读数为 $0.90\UV$；若将电阻箱拨到如图所示\tkanswer{$110\UO$}时，伏特表的读数如右图所示是\tkanswer{$1.10\UV$}。

	\item
	根据以上实验数据，可以算出该节干电池的电动势 $\varepsilon=$\tkanswer{$1.3\UV$}，内电阻 $r=$\tkanswer{$20\UO$}。
\end{enumerate}
\twopicture[align=right, showlabel=false, h=3cm]{figs/22a.png}{figs/22b.png}

%% answer: （1）见解析；（2）110 Ω，1.10 V；（3）1.3 V，20 Ω
\jdanswer{
\begin{enumerate}
	\item
	正确画出电路图：将电源、电键和电阻箱串联，伏特表并联在电阻箱的两端。

	\item
	$110\UO$，$1.10\UV$

	\item
	$1.3\UV$，$20\UO$
\end{enumerate}
}

%% memo:
\memoanswer{本题原卷及材料中未提供详解，待补充。}

\item
%% number: 23
%% paperName: 1986年上海高考第 23 题 5 分
%% typeId: 4
%% type: 实验题
%% chapter: 电路
%% point: 故障分析
%% method: 无
%% score: 5
%% degree: 450
%% duplicateId: 0
%% body:
右图中，$L_{1}$、$L_{2}$ 是两个相同的乳白色的小灯泡，额定电压均为 $3\UV$。$E$ 是一电源，电动势为 $6\UV$。电路接通后灯泡不亮，而导线和各连接点的接触均良好。给你一个伏特表（量程 $0\sim3\sim15\UV$），要求在不断开电路的情况下，查出是哪一个器件已损坏。试述操作步骤（不讨论有两个或三个器件同时损坏的情况）。
\onepicture[align=right, w=4cm]{\includesvg{figs/23}}

%% answer: 见解析
\jdanswer{
\begin{steps}
	\item
	伏特表量程选用 $15\UV$ 档。

	\item
	将伏特表并联在电池两端，若读数显著低于 $6\UV$，则电池不能用；若在 $6\UV$ 左右，则电池正常，再查灯泡。

	\item
	用伏特表测某一个灯泡的电压，若读数为零，则另一灯泡的灯丝已断；若读数为 $6\UV$ 左右，则这个灯泡的灯丝断。
\end{steps}
}

%% memo:
\memoanswer{本题原卷及材料中未提供详解，待补充。}

\item
%% number: 24
%% paperName: 1986年上海高考第 24 题 6 分
%% typeId: 6
%% type: 计算题
%% chapter: 曲线运动
%% point: 平抛运动
%% method: 无
%% score: 6
%% degree: 650
%% duplicateId: 0
%% body:
以初速 $v_{0}$ 水平抛出一物体，试求：
\begin{enumerate}
	\item
	抛出多少时间后，物体的竖直位移和水平位移相等？

	\item
	当竖直速度与水平速度相等时，物体的竖直位移和水平位移之比是多少？
\end{enumerate}

%% answer: （1）\frac{2v_{0}}{g}；（2）1:2
\jdanswer{
\begin{enumerate}
	\item
	$\frac{2v_{0}}{g}$

	\item
	$1:2$
\end{enumerate}
}

%% memo:
\memoanswer{
（1）$x=v_{0}t$，$y=\frac{1}{2}gt^{2}$，由 $x=y$ 得

$v_{0}t=\frac{1}{2}gt^{2}$，

得出：$t=\frac{2v_{0}}{g}$。

（2）$v_{y}=gT=v_{0}$，$T=\frac{v_{0}}{g}$；

$\frac{y}{x}=\frac{\frac{1}{2}gT^{2}}{v_{0}T}=\frac{gT}{2v_{0}}$，把 $T=\frac{v_{0}}{g}$ 代入得

$\frac{y}{x}=\frac{1}{2}$。
}

\item
%% number: 25
%% paperName: 1986年上海高考第 25 题 6 分
%% typeId: 6
%% type: 计算题
%% chapter: 电磁感应
%% point: 电磁感应能量分析
%% method: 无
%% score: 6
%% degree: 650
%% duplicateId: 0
%% body:
一边长 $l=0.50\Um$、电阻 $R=0.40\UO$ 的正方形线圈，在外力作用下，以 $v=2.0\Ums$ 的速度匀速进入磁感应强度 $B=0.04\UT$ 的匀强磁场，线圈平面与磁感线垂直。试求线圈在图示位置时外力的功率。
\onepicture[align=right, w=5.5cm]{\includesvg{figs/25}}

%% answer: 4\times10^{-3} W
\jdanswer{
$4\times10^{-3}\UW$
}

%% memo:
\memoanswer{
$E=Blv$

$I=\frac{E}{R}$

$F=BIl$

维持线圈匀速运动外力大小 $F'=F$

$P=F'v$

$=\frac{B^{2}l^{2}v^{2}}{R}=\frac{0.04^{2}\times0.50^{2}\times2.0^{2}}{0.40}\UW=4\times10^{-3}\UW$
}

\item
%% number: 26
%% paperName: 1986年上海高考第 26 题 11 分
%% typeId: 6
%% type: 计算题
%% chapter: 磁场
%% point: 带电粒子在磁场中的运动
%% method: 无
%% score: 11
%% degree: 650
%% duplicateId: 0
%% body:
一平行板电容器的两极板 M、N 通过变阻器 $R$ 与电源 $E$ 相连。在距 N 板较远处，有一磁感应强度 $B$ 为 $0.10\UT$ 的匀强磁场，方向垂直纸面向外，磁场的边界面 P 与 N 板平行。不考虑重力场。

当极板间的电压为 $4.5\UV$ 时，一质量为 $1.0\times10^{-12}\Ukg$、电量为 $1.0\times10^{-8}\UC$ 的点电荷 $q$，从 M 板由静止开始加速，并从 N 板上的小孔 C 射出，经金属屏蔽管 G 进入磁场。
\begin{enumerate}
	\item
	若要 $q$ 经磁场后从 N 板上的小孔 D 射回电容器内，则 C、D 之间的距离应是多少？

	\item
	当 $q$ 回到 M 板时的速度是多大？
\end{enumerate}

若在磁场区域内放置另一带负电的点电荷 $Q$，调节极板间的电压，再使 $q$ 从 M 板出发经 C 孔和 G 管进入磁场后，在磁场区域内仍沿第 1 问中的运动轨迹运动，则
\begin{enumerate}
	\item
	点电荷 $Q$ 应放在哪里？

	\item
	变阻器 $R$ 的滑键应向哪边移动？为什么？
\end{enumerate}
\onepicture[align=right, w=6cm]{\includesvg{figs/26}}

%% answer: （1）0.6 m；（2）0；（3）圆轨道的圆心处；（4）向右端移动
\jdanswer{
\begin{enumerate}
	\item
	$CD=0.6\Um$

	\item
	$v=0$

	\item
	点电荷 $Q$ 应放在磁场中圆轨道的圆心处。

	\item
	变阻器的滑键向右端移动。理由：a）点电荷 $q$ 在洛仑兹力和库仑力作用下作圆周运动，向心力增大；b）保持原轨道运动，$v$ 要增大；c）电容器两板电压 $U$ 要加大，所以变阻器的滑键向右移。
\end{enumerate}
}

%% memo:
\memoanswer{
（1）$\frac{1}{2}mv^{2}=qU$

$v=\sqrt{\frac{2qU}{m}}=\sqrt{\frac{2\times10^{-8}\times4.5}{10^{-12}}}\Ums=300\Ums$

$qvB=\frac{mv^{2}}{R}$

$CD=2R=\frac{2mv}{qB}=\frac{2\times10^{-12}\times300}{10^{-8}\times0.10}\Um=0.6\Um$

（2）$v=0$

①点电荷 $Q$ 应放在磁场中圆轨道的圆心处。

②变阻器的滑键向右端移动。

理由：a）点电荷 $q$ 在洛仑兹力和库仑力作用下作圆周运动，向心力增大。b）保持原轨道运动，$v$ 要增大。c）电容器两板电压 $U$ 要加大，所以变阻器的滑键向右移。
}

\item
%% number: 27
%% paperName: 1986年上海高考第 27 题 9 分
%% typeId: 6
%% type: 计算题
%% chapter: 气体性质
%% point: 玻意耳定律
%% method: 无
%% score: 9
%% degree: 450
%% duplicateId: 0
%% body:
如图所示，可沿缸壁自由滑动的活塞将圆筒形气缸分隔成 A、B 两部分，气缸的底部通过装有阀门 K 的细管与一密闭容器 C 相连。活塞与气缸的顶部间连有一弹簧，当活塞位于气缸底部时，弹簧恰好无形变。开始时，B 内充有一定量的气体，A、C 内为真空。B 部分的高 $l_{1}=0.10\Um$，B 与 C 的容积正好相等。此时弹簧对活塞的作用力正好等于活塞的重力。现将阀门打开，并将整个装置倒置。当达到新的平衡时，B 部分的高 $l_{2}$ 等于多少米？
\onepicture[align=right, h=5.5cm]{\includesvg{figs/27}}

%% answer: 0.17 m
\jdanswer{
$0.17\Um$
}

%% memo:
\memoanswer{
设活塞质量为 $m$，弹簧倔强系数为 $k$，气缸的横截面面积为 $S$，当 B 室内气体压强为 $p_{0}$ 时，弹簧伸长 $l_{1}$，

$mg=kl$；

$p_{0}S=mg+kl_{1}$

$p_{0}=\frac{2kl_{1}}{S}$ 或 $p_{0}=\frac{2mg}{S}$（1）

设开启阀门将整个装置又倒置后，B 容器内气体压强为 $p$，B 部分的高为 $l_{2}$，

$p_{0}l_{1}S=p(l_{1}+l_{2})S$

$p=\frac{l_{1}}{l_{1}+l_{2}}p_{0}=\frac{2kl_{1}^{2}}{(l_{1}+l_{2})S}$（2）

$pS=kl_{2}-mg$

$p=\frac{k(l_{2}-l_{1})}{S}$（3）

由（2）、（3）式得：$l_{2}=\sqrt{3}l_{1}=\sqrt{3}\times0.1\Um=0.17\Um$。
}

\end{enumerate}

\end{document}
